\section{A representation-level dichotomy}
\label{sec:dichotomy}

The coding-tree obstruction by itself does not distinguish a difficult PDE from a difficult representation. The following benchmark shows that the latter distinction is real.

Fix $\eta\neq0$ and
\begin{equation*}
  0<a\leq
  \frac{\erfc(1/\sqrt2)}{\sqrt3},
  \qquad
  \erfc(r)=\frac2{\sqrt\pi}\int_r^\infty e^{-s^2}\,ds.
\end{equation*}
Consider
\begin{equation}
  \partial_tu+\frac12u_{xx}
  +\eta\pa{e^{(u_x)^4}-1}=0,
  \qquad
  u(T,x)=a\cos x.
  \label{eq:dichotomy-benchmark}
\end{equation}
Put
\begin{equation}
  T_*(\eta)
  =
  \frac1{2\pi e(e-1)^2\eta^2}.
  \label{eq:T-star}
\end{equation}

For the marked branching representation of Henry-Labord\`ere, Oudjane, Tan, Touzi, and Warin~\cite{HenryLabordereEtAl2019}, take monomial types and coefficients
\begin{equation*}
  L=\set{(0,4r):r\geq1},
  \qquad
  c_r=\frac\eta{r!},
  \qquad
  p_r=\frac1{(e-1)r!},
\end{equation*}
and lifetime density
\begin{equation}
  \rho_T(s)
  =
  \frac{s^{-1/2}e^{-s/(2T)}}{\sqrt{2\pi T}},
  \qquad s>0.
  \label{eq:HLOTW-lifetime}
\end{equation}

\begin{theorem}[Representation-level dichotomy]
\label{thm:representation-dichotomy}
For~\eqref{eq:dichotomy-benchmark} the following statements hold.

\begin{enumerate}[label=\textup{(\roman*)}]
\item Let
\begin{equation*}
  f(z_0,z_1)=\eta(e^{z_1^4}-1).
\end{equation*}
For every $T>0$, $0\leq t<T$, and $x\in\R$, the Nguwi--Penent--Privault coding-tree functionals satisfy
\begin{equation*}
  \E\abs{H(\mathcal T_{t,x,f^*})}=\infty,
  \qquad
  \E\abs{H(\mathcal T_{t,x,\Id})}=\infty.
\end{equation*}
The conclusion holds for every strictly positive coding-tree lifetime density and the positive mechanism-selection probabilities of the construction.

\item If $0<T<T_*(\eta)$ and the HLOTW estimator uses the data above, then its marked branching functional $\psi_{t,x}$ satisfies
\begin{equation*}
  \E\abs{\psi_{t,x}}^2\leq1,
  \qquad
  0\leq t\leq T,
  \quad x\in\R.
\end{equation*}
Moreover
\begin{equation*}
  u(t,x)=\E[\psi_{t,x}]
\end{equation*}
is a continuous viscosity solution of~\eqref{eq:dichotomy-benchmark}.
\end{enumerate}
\end{theorem}

The lifetime density~\eqref{eq:HLOTW-lifetime} is singular at the origin. The theorem uses the positive-time extension of the HLOTW moment argument: internal lifetimes are almost surely positive, and the proof of the second-moment bound evaluates $\rho$ only at positive arguments. This is the only modification to the literal endpoint-continuity wording of the published assumptions.

\proofplaceholder{Insert the audited proof. The NPP half should be deduced from the repeated-Hessian obstruction using the $(4r)!/r!$ lower bound. The HLOTW half should verify the explicit $q=2$ majorant, including the singular-lifetime endpoint extension, without presenting the dichotomy as the main theme of the manuscript.}
